% !TEX root = ../main.tex
% ============================================================================
%  Estilo
% ============================================================================

\definecolor{corprincipal}{HTML}{\corhex}
\definecolor{fundocodigo}{gray}{0.96}
\colorlet{azulclaro}{corprincipal!35!white}

\newcommand{\email}[1]{\href{mailto:#1}{\nolinkurl{#1}}}

\onehalfspacing
\setlength{\parindent}{1.5em}
\setlist{itemsep=0.2em}

\hypersetup{
  colorlinks = true,
  linkcolor  = black,
  citecolor  = corprincipal,
  urlcolor   = corprincipal,
  pdftitle   = {\titulo},
}

% --- Nomes em português -----------------------------------------------------
\addto\captionsportuguese{%
  \renewcommand{\contentsname}{Índice}%
  \renewcommand{\listfigurename}{Índice de Figuras}%
  \renewcommand{\listtablename}{Índice de Tabelas}%
  \renewcommand{\appendixname}{Anexo}%
}
\renewcommand{\lstlistingname}{Código}
\renewcommand{\lstlistlistingname}{Índice de Código}

\newcommand{\crefrangeconjunction}{ a\nobreakspace}
\newcommand{\crefpairconjunction}{ e\nobreakspace}
\newcommand{\crefmiddleconjunction}{, }
\newcommand{\creflastconjunction}{ e\nobreakspace}
\crefname{figure}{figura}{figuras}        \Crefname{figure}{Figura}{Figuras}
\crefname{table}{tabela}{tabelas}          \Crefname{table}{Tabela}{Tabelas}
\crefname{equation}{equação}{equações}     \Crefname{equation}{Equação}{Equações}
\crefname{chapter}{capítulo}{capítulos}    \Crefname{chapter}{Capítulo}{Capítulos}
\crefname{section}{secção}{secções}        \Crefname{section}{Secção}{Secções}
\crefname{subsection}{secção}{secções}     \Crefname{subsection}{Secção}{Secções}
\crefname{appendix}{anexo}{anexos}         \Crefname{appendix}{Anexo}{Anexos}
\crefname{listing}{código}{códigos}        \Crefname{listing}{Código}{Códigos}
\crefname{lstlisting}{código}{códigos}     \Crefname{lstlisting}{Código}{Códigos}

% --- Títulos ----------------------------------------------------------------
\titleformat{\chapter}[hang]
  {\normalfont\huge\bfseries\color{corprincipal}}{\thechapter}{0.6em}{}
\titlespacing*{\chapter}{0pt}{0pt}{1.2em}
\titleformat{\section}
  {\normalfont\Large\bfseries\color{corprincipal}}{\thesection}{0.8em}{}
\titleformat{\subsection}
  {\normalfont\large\bfseries\color{corprincipal}}{\thesubsection}{0.8em}{}
\titleformat{\subsubsection}
  {\normalfont\normalsize\bfseries\color{corprincipal}}{\thesubsubsection}{0.8em}{}

% --- Legendas ---------------------------------------------------------------
\captionsetup{labelfont={bf, color=corprincipal}, font=small, labelsep=period}

% --- Índice -----------------------------------------------------------------
\renewcommand{\cftchapleader}{\cftdotfill{\cftdotsep}}
\setlength{\cftbeforechapskip}{0.4em}

% --- Cabeçalhos e rodapés ---------------------------------------------------
\fancypagestyle{plain}{%            páginas de início de capítulo / pré-textuais
  \fancyhf{}%
  \fancyfoot[C]{\thepage}%
  \renewcommand{\headrulewidth}{0pt}%
}
\fancypagestyle{corpo}{%
  \fancyhf{}%
  \fancyhead[L]{\small\color{corprincipal}\nouppercase{\leftmark}}%
  \fancyhead[R]{\small\color{corprincipal}\titulocurto}%
  \fancyfoot[C]{\thepage}%
  \renewcommand{\headrulewidth}{0.4pt}%
}
\renewcommand{\chaptermark}[1]{\markboth{\thechapter.\ #1}{}}

% --- Código -----------------------------------------------------------------
\lstset{
  backgroundcolor   = \color{fundocodigo},
  basicstyle        = \ttfamily\footnotesize,
  keywordstyle      = \color{corprincipal}\bfseries,
  commentstyle      = \color{gray}\itshape,
  stringstyle       = \color{teal!70!black},
  numbers           = left,
  numberstyle       = \tiny\color{gray},
  frame             = single,
  rulecolor         = \color{gray!40},
  breaklines        = true,
  tabsize           = 2,
  showstringspaces  = false,
  captionpos        = b,
  % listings não suporta UTF-8: sem isto, acentos dentro de código dão erro
  literate =
    {á}{{\'a}}1 {à}{{\`a}}1 {â}{{\^a}}1 {ã}{{\~a}}1
    {é}{{\'e}}1 {ê}{{\^e}}1 {í}{{\'i}}1
    {ó}{{\'o}}1 {ô}{{\^o}}1 {õ}{{\~o}}1 {ú}{{\'u}}1 {ü}{{\"u}}1
    {Á}{{\'A}}1 {À}{{\`A}}1 {Â}{{\^A}}1 {Ã}{{\~A}}1
    {É}{{\'E}}1 {Ê}{{\^E}}1 {Í}{{\'I}}1
    {Ó}{{\'O}}1 {Ô}{{\^O}}1 {Õ}{{\~O}}1 {Ú}{{\'U}}1
    {ç}{{\c c}}1 {Ç}{{\c C}}1 {º}{{\textordmasculine}}1 {ª}{{\textordfeminine}}1
    {€}{{\texteuro}}1 {°}{{\textdegree}}1,
}

% --- Acrónimos --------------------------------------------------------------
\newcommand{\printacronimos}{%
  \printnoidxglossary[type=\acronymtype, title={Lista de Acrónimos}, toctitle={Lista de Acrónimos}]}

% --- Transições de numeração ------------------------------------------------
\newcommand{\iniciarcorpo}{%
  \clearpage
  \pagenumbering{arabic}%
  \pagestyle{corpo}%
}

\newcommand{\iniciaranexos}{%
  \appendix
  \titleformat{\chapter}[hang]
    {\normalfont\huge\bfseries\color{corprincipal}}{\appendixname~\thechapter}{0.6em}{}%
  \renewcommand{\chaptermark}[1]{\markboth{\appendixname~\thechapter.\ ##1}{}}%
  \addtocontents{toc}{%
    \protect\renewcommand{\protect\cftchappresnum}{\appendixname~}%
    \protect\setlength{\protect\cftchapnumwidth}{5.2em}}%
  \renewcommand{\thepage}{\thechapter-\arabic{page}}%
}

% Usar \anexo{Título} em vez de \chapter nos anexos: reinicia a página (A-1, B-1, ...)
\newcommand{\anexo}[1]{\chapter{#1}\setcounter{page}{1}}
